\section{Evaluation Metrics and Statistical Protocol}
\label{app:evaluation_metrics}

We evaluate the Q7 reasoning response (Seg2) on the fixed cohort of 162
participants. The proposed method returns one ordered array of aligned units,
where each unit contains a C--K state and the text realizing that state. The
state alphabet is
\(\mathcal Z=\{K,C,K{\to}C,C{\to}K,K{\to}K,C{\to}C\}\).
V30 denotes the participant's value profile and is not a trajectory state.
Let \(M_u=(m_{u1},\ldots,m_{uT_u})\) and
\(\widehat M_u=(\widehat m_{u1},\ldots,\widehat m_{u\widehat T_u})\)
be the gold and predicted trajectories for participant \(u\), respectively.
Auxiliary boundary symbols, if present, are removed before evaluation. Let
\(Y_u\) and \(\widehat Y_u\) denote the corresponding gold and generated
responses. Metrics are grouped into trajectory fidelity, participant-specific
personalization, and anti-collapse diagnostics; high diversity alone is not
interpreted as fidelity or personalization.

\subsection{Textual Own-User Fidelity}

\paragraph{BGE own-user margin.}
We encode every response with the frozen BGE-M3 dense encoder \(e(\cdot)\)
(1024 dimensions; L2-normalized). For
\[
q_{uv}=\cos\!\left(e(\widehat Y_u),e(Y_v)\right),
\]
the participant-level margin is
\begin{equation}
b_u=q_{uu}-\frac{1}{N-1}\sum_{v\ne u}q_{uv},
\qquad
\operatorname{BGE\text{-}Margin}=\frac{1}{N}\sum_u b_u .
\label{eq:bge_own_user_margin}
\end{equation}
A positive value means that a generated response is, on average, closer to
its intended participant's gold response than to the gold responses of the
other participants. This is an identity-specific text metric, rather than a
generic semantic-similarity score. We also retain mean own-user and
other-user similarities for interpretation.

\subsection{C--K Trajectory Fidelity}

\paragraph{Exact match and length error.}
Exact match requires equality of the complete ordered state sequences:
\begin{equation}
\operatorname{EM}=\frac{1}{N}\sum_u
\mathbb I[\widehat M_u=M_u].
\label{eq:kc_exact_match}
\end{equation}
Length error measures only the discrepancy in the number of operations:
\begin{equation}
\operatorname{LenErr}=\frac{1}{N}\sum_u
\left|\widehat T_u-T_u\right|.
\label{eq:kc_length_error}
\end{equation}
Thus, EM is higher-is-better and LenErr is lower-is-better.

\paragraph{KC normalized edit similarity (KC-NES).}
Let \(\operatorname{Lev}(\cdot,\cdot)\) be Levenshtein distance with unit
insertion, deletion, and substitution costs. We compute
\begin{equation}
\operatorname{NES}_u
=1-\frac{\operatorname{Lev}(\widehat M_u,M_u)}
{\max(\widehat T_u,T_u,1)},
\qquad
\operatorname{KC\text{-}NES}=\frac{1}{N}\sum_u\operatorname{NES}_u .
\label{eq:kc_nes}
\end{equation}
KC-NES lies in \([0,1]\); higher values indicate better joint agreement in
state identity, order, and length. An invalid or missing trajectory receives
zero under the fixed failure policy.

\subsection{C--K Trajectory Personalization}

KC-NES measures fidelity to the intended gold trajectory but does not by
itself establish that the prediction is more specific to its owner than to a
population-level trajectory. We therefore contrast each prediction with all
gold trajectories in the cohort. Define
\begin{equation}
a_{uv}=1-\frac{\operatorname{Lev}(\widehat M_u,M_v)}
{\max(\widehat T_u,T_v,1)}.
\label{eq:kc_cross_user_similarity}
\end{equation}

\paragraph{KC own-user margin.}
\begin{equation}
\Delta_u^{\mathrm{KC}}
=a_{uu}-\frac{1}{N-1}\sum_{v\ne u}a_{uv},
\qquad
\operatorname{KC\text{-}Margin}
=\frac{1}{N}\sum_u\Delta_u^{\mathrm{KC}} .
\label{eq:kc_own_user_margin}
\end{equation}
A positive margin indicates that the generated trajectory fits its intended
participant better than an average non-target participant.

\paragraph{KC own-user AUC.}
For each participant, we treat the own gold trajectory as the positive and
each other participant's trajectory as a negative:
\begin{equation}
\operatorname{AUC}_u^{\mathrm{KC}}
=\frac{1}{N-1}\sum_{v\ne u}
\left(
\mathbb I[a_{uu}>a_{uv}]
+\frac{1}{2}\mathbb I[a_{uu}=a_{uv}]
\right),
\qquad
\operatorname{KC\text{-}AUC}
=\frac{1}{N}\sum_u\operatorname{AUC}_u^{\mathrm{KC}} .
\label{eq:kc_own_user_auc}
\end{equation}
KC-AUC is bounded by \([0,1]\); \(0.5\) is the chance-ranking reference.
KC-AUC and KC-Margin are the primary sequence-personalization metrics.

\subsection{Trajectory Anti-Collapse and Geometry}

Let
\[
d(A,B)=1-\left[1-
\frac{\operatorname{Lev}(A,B)}{\max(|A|,|B|,1)}\right]
\]
be normalized edit distance. Let
\(\overline d_{\mathrm{pred}}\) and \(\overline d_{\mathrm{gold}}\)
be its means over all unordered participant pairs for predicted and gold
trajectories, respectively.

\paragraph{KC diversity retention (KCDR).}
\begin{equation}
\operatorname{KCDR}
=\frac{\overline d_{\mathrm{pred}}}{\overline d_{\mathrm{gold}}}.
\label{eq:kcdr}
\end{equation}
Higher KCDR indicates greater retention of the cross-participant variation
present in the gold trajectories, with \(1\) denoting equal average
dispersion. Values above \(1\) indicate over-dispersion and therefore are not
automatically evidence of better personalization.

\paragraph{Trajectory geometry correlation (TGC).}
\begin{equation}
\operatorname{TGC}
=\rho_{\mathrm{Spearman}}
\left(
\{d(\widehat M_u,\widehat M_v)\}_{u<v},
\{d(M_u,M_v)\}_{u<v}
\right).
\label{eq:tgc}
\end{equation}
TGC lies in \([-1,1]\) and tests whether the relative geometry of
between-participant trajectory distances is preserved; higher is better.

\paragraph{Unique trajectories and top share.}
For each distinct predicted trajectory \(A\), let
\(n_A=\sum_u\mathbb I[\widehat M_u=A]\). We report
\begin{equation}
\operatorname{Unique}=\left|\{A:n_A>0\}\right|,
\qquad
\operatorname{TopShare}=\frac{\max_A n_A}{N}.
\label{eq:unique_topshare}
\end{equation}
More unique trajectories and a lower TopShare indicate less exact mode
collapse. Chain diversity equals \(1-\operatorname{TopShare}\) and is
therefore omitted from the main table as redundant. KCDR, TGC, Unique, and
TopShare are cohort-level diagnostics: none alone establishes that the
variation is assigned to the correct participants.

\subsection{NLI-Based V30 Value Expression}

A frozen two-pole NLI probe maps each generated response and each human gold
response to one stance for every V30 dimension:
\(\{\textsc{agree},\textsc{disagree},\textsc{not-expressed}\}\).
Accordingly, ``Human'' below denotes automatic agreement with value claims
extracted from the human response, not a manual rating. For participant \(u\),
let \(G_u\) and \(H_u\) be the sets of non-neutral
\((\text{dimension},\text{polarity})\) claims extracted from generated and
human responses.

\paragraph{Human expression precision, recall, and F1.}
We micro-aggregate directional claims over participants:
\begin{align}
\operatorname{TP}&=\sum_u|G_u\cap H_u|, &
\operatorname{FP}&=\sum_u|G_u\setminus H_u|, &
\operatorname{FN}&=\sum_u|H_u\setminus G_u|, \nonumber\\
P_{\mathrm{human}}&=\frac{\operatorname{TP}}
{\operatorname{TP}+\operatorname{FP}},&
R_{\mathrm{human}}&=\frac{\operatorname{TP}}
{\operatorname{TP}+\operatorname{FN}},&
F1_{\mathrm{human}}&=\frac{2P_{\mathrm{human}}R_{\mathrm{human}}}
{P_{\mathrm{human}}+R_{\mathrm{human}}}.
\label{eq:human_expression_prf}
\end{align}
All three are higher-is-better. Matching a dimension with the wrong polarity
does not count as a true positive. Undefined denominators are reported as
missing rather than assigned an artificial perfect score.

\paragraph{Profile-polarity precision and expression coverage.}
The six-point V30 score is binarized as \textsc{disagree} for scores 1--3 and
\textsc{agree} for scores 4--6. Let \(p_{uj}\) be this profile polarity and
\(g_{uj}\) the generated claim polarity. For participants with
\(|G_u|>0\),
\begin{equation}
\operatorname{PPP}_u
=\frac{1}{|G_u|}\sum_{(j,g_{uj})\in G_u}\mathbb I[g_{uj}=p_{uj}].
\label{eq:profile_polarity_precision}
\end{equation}
Profile-polarity precision is the mean of \(\operatorname{PPP}_u\) over
participants who express at least one claim. Expression coverage includes
all participants and is
\begin{equation}
\operatorname{Coverage}=\frac{1}{N}\sum_u\frac{|G_u|}{30}.
\label{eq:expression_coverage}
\end{equation}
PPP measures directional consistency conditional on expression, whereas
coverage exposes systems that attain precision by expressing very little.

\paragraph{Own-profile advantage.}
Let
\[
\operatorname{Fit}(G_u,p_v)
=\frac{1}{|G_u|}\sum_{(j,g_{uj})\in G_u}\mathbb I[g_{uj}=p_{vj}]
\]
for \(|G_u|>0\). The participant-specific contrast is
\begin{equation}
\operatorname{OPA}_u
=\operatorname{Fit}(G_u,p_u)
-\frac{1}{N-1}\sum_{v\ne u}\operatorname{Fit}(G_u,p_v),
\qquad
\operatorname{OPA}=\operatorname{mean}_{u:|G_u|>0}\operatorname{OPA}_u .
\label{eq:own_profile_advantage}
\end{equation}
A positive OPA means that the expressed value polarities fit the target
profile better than the profiles of other participants on average. As with
PPP, participants without generated claims are excluded, so OPA must be
interpreted jointly with coverage.

\subsection{Statistical Reporting}

The participant is the sole unit of inference. All system comparisons use a
paired participant bootstrap with 10,000 resamples; each resample draws 162
participants with replacement, recomputes the statistic for both systems,
and reports the 2.5th and 97.5th percentiles of the paired difference as its
95\% confidence interval. Participant-decomposable scores are averaged within
each resample, while micro-aggregated value-expression scores are recomputed
from the resampled claim counts. For KCDR, TGC, Unique, and TopShare, the
complete cohort statistic is recomputed after participant resampling.
Although KCDR and TGC contain participant pairs, these pairs are dependent
and are never treated as independent observations. We report valid sample
sizes and expression coverage wherever undefined or missing outputs can
change the effective cohort.
